\documentclass[12pt]{article}

\usepackage{amsmath, amssymb}
\usepackage{geometry}
\usepackage{graphicx}
\usepackage{setspace}

\geometry{margin=1in}

\title{Lecture Scribe \\ \large CSE400: Fundamentals of Probability in Computing}
\author{Lecture L15--L18 \\ Joint Random Variables and Joint Distributions}
\date{Instructor: Dr. Dhaval Patel \\ March 10, 2026}

\begin{document}

\maketitle

\onehalfspacing

\section{Introduction}

This lecture introduces the concept of \textbf{Joint Random Variables} and their associated probability distributions. 
In many real-world applications, multiple random variables must be analyzed simultaneously because they may influence one another.

Instead of studying each variable independently, joint distributions allow us to understand the \textbf{relationship between two random variables}.

\section{Why Study Joint Random Variables?}

In practical systems, multiple uncertain quantities often interact with each other. 
To analyze such systems, we must study the joint behavior of random variables.

Several motivating examples were discussed.

\subsection{Example 1: Smart Transportation}

Consider a smart traffic system with two random variables:

\[
X = \text{Vehicle Speed}
\]

\[
Y = \text{Traffic Density}
\]

Question:

\begin{quote}
Can vehicle speed be high when traffic density is also high?
\end{quote}

Since traffic density may affect vehicle speed, the variables are \textbf{correlated}. 
Therefore, studying their joint distribution allows probabilistic analysis for applications such as:

\begin{itemize}
\item Adaptive traffic signals
\item Congestion prediction
\item Autonomous vehicle routing
\end{itemize}

\subsection{Example 2: Wireless Network Performance}

In a wireless communication system:

\[
X = \text{Signal-to-Noise Ratio (SNR)}
\]

\[
Y = \text{Data Throughput}
\]

Throughput depends on SNR, but other factors such as interference and congestion also influence performance.

Question:

\begin{quote}
If SNR is high, will throughput always be high?
\end{quote}

The answer is not always, because network congestion can affect throughput. 
Thus, analyzing both variables jointly is necessary.

Applications include:

\begin{itemize}
\item 5G scheduling
\item Adaptive modulation
\item Network planning
\end{itemize}

\subsection{Example 3: Weather Prediction}

Two random variables:

\[
X = \text{Rainfall}
\]

\[
Y = \text{Temperature}
\]

Weather forecasting requires studying how these variables vary together.

\subsection{Example 4: Drone Delivery}

Consider:

\[
X = \text{Wind Speed}
\]

\[
Y = \text{Delivery Time}
\]

Higher wind speeds may increase delivery time, making joint analysis necessary for UAV logistics systems.

\subsection{Example 5: Rolling Two Dice}

Let

\[
X = \text{value of first die}
\]

\[
Y = \text{value of second die}
\]

Questions such as:

\[
X + Y = 7
\]

or

\[
X > Y
\]

require computing probabilities involving both variables simultaneously, which is done using a \textbf{Joint Probability Mass Function}.

\section{Discrete Case: Joint Probability Mass Function}

When the random variables are discrete, their joint behavior is described by a \textbf{Joint PMF}.

\subsection{Definition}

For discrete random variables \(X\) and \(Y\), the joint PMF is defined as

\[
p_{X,Y}(x,y) = P(X=x, Y=y)
\]

This represents the probability that:

\begin{itemize}
\item \(X\) takes value \(x\)
\item \(Y\) takes value \(y\)
\end{itemize}

simultaneously.

\subsection{Joint PMF Table Representation}

A joint PMF is often represented using a table where:

\begin{itemize}
\item Rows correspond to values of \(X\)
\item Columns correspond to values of \(Y\)
\item Each entry contains \(P(X=x, Y=y)\)
\end{itemize}

The sum of all probabilities in the table must equal:

\[
\sum_x \sum_y p_{X,Y}(x,y) = 1
\]

\section{Example: Rolling Two Dice}

Consider rolling two fair dice.

Let

\[
X = \text{outcome of the first die}
\]

\[
Y = \text{outcome of the second die}
\]

Each die can take values from 1 to 6.

Since the dice are fair, every pair \((x,y)\) has equal probability.

Total outcomes:

\[
6 \times 6 = 36
\]

Thus,

\[
P(X=x, Y=y) = \frac{1}{36}
\]

for all valid pairs \((x,y)\).

Using the joint PMF table, probabilities of events such as

\[
P(X+Y=7)
\]

can be computed by summing all pairs satisfying the condition.

\section{Continuous Case: Joint Probability Density Function}

When random variables are continuous, the joint distribution is described using a \textbf{Joint Probability Density Function (PDF)}.

\subsection{Definition}

The joint PDF of continuous random variables \(X\) and \(Y\) is denoted by

\[
f_{X,Y}(x,y)
\]

Probabilities are computed using integration over a region.

\[
P((X,Y) \in R) = \int\int_R f_{X,Y}(x,y)\,dx\,dy
\]

\subsection{Geometric Interpretation}

For continuous variables, probabilities correspond to the \textbf{volume under the joint PDF surface} over a specified region.

Thus, probability is determined by integrating the density function over the desired area.

\section{Worked Example (Continuous Case)}

To compute probabilities in the continuous case:

\begin{enumerate}
\item Identify the joint PDF \(f_{X,Y}(x,y)\)
\item Determine the region \(R\) corresponding to the event
\item Evaluate the double integral
\end{enumerate}

\[
P((X,Y)\in R) = \int\int_R f_{X,Y}(x,y)\,dx\,dy
\]

The example solution in the lecture demonstrates computing probability by integrating over the specified region.

\section{Joint Cumulative Distribution Function (Joint CDF)}

Another way to describe the joint behavior of two random variables is through the \textbf{Joint CDF}.

\subsection{Definition}

The joint CDF of \(X\) and \(Y\) is defined as

\[
F_{X,Y}(x,y) = P(X \le x, Y \le y)
\]

This function gives the probability that:

\begin{itemize}
\item \(X\) is less than or equal to \(x\)
\item \(Y\) is less than or equal to \(y\)
\end{itemize}

simultaneously.

\subsection{Continuous Case}

For continuous random variables, the joint CDF can be computed from the joint PDF as

\[
F_{X,Y}(x,y) =
\int_{-\infty}^{x}
\int_{-\infty}^{y}
f_{X,Y}(u,v)\,dv\,du
\]

\subsection{Discrete Case}

For discrete variables, the joint CDF is computed using summation:

\[
F_{X,Y}(x,y) =
\sum_{u \le x}
\sum_{v \le y}
P(X=u, Y=v)
\]

\section{Properties of the Joint CDF}

Important properties include:

\begin{itemize}
\item \(0 \le F_{X,Y}(x,y) \le 1\)
\item \(F_{X,Y}(x,y)\) is non-decreasing in both \(x\) and \(y\)
\item As \(x\) and \(y\) approach infinity, the CDF approaches 1
\end{itemize}

\section{Summary}

This lecture introduced the concept of analyzing multiple random variables simultaneously.

Key concepts covered:

\begin{itemize}
\item Motivation for studying joint random variables
\item Joint PMF for discrete variables
\item Joint PDF for continuous variables
\item Geometric interpretation of joint density
\item Joint cumulative distribution function
\end{itemize}

Joint distributions allow us to analyze relationships between variables and are essential for modeling real-world probabilistic systems.

\end{document}